\documentclass[spanish, 12pt, a4paper]{article}

\usepackage{name}

\begin{document}

\maketitle
\thispagestyle{fancy}
\setcounter{page}{1}
\maketitle
\begin{multicols}{2}

\section{El mapeo de Beverton - Holt}
El modelo de Beverton-Holt es un modelo de dinámica de poblaciones que reproduce el comportamiento de la ecuación logística, y está definido por la siguiente expresión:

\begin{equation}
n_{t+1} = \frac{r n_t}{1 + \frac{r-1}{K} n_t},
\label{eqn:map-beverton-holt}
\end{equation}

donde $n_t$ es la población discreta en el tiempo $t$, $K$ es la capacidad de carga y $r$ es la tasa de reproducción.

En primer lugar, buscamos los puntos de equilibrio, es decir, aquellos valores de $n_t$ para los cuales $n_{t+1} = n_t$ en la ecuación \eqref{eqn:map-beverton-holt}. De esta manera, obtenemos $n_t = K$ y $n_t = 0$, al igual que en el modelo logístico. Para analizar la estabilidad de estos puntos de equilibrio, calculamos la derivada de la ecuación \eqref{eqn:map-beverton-holt} con respecto a $n_t$ y la evaluamos en los puntos de equilibrio. Así, obtenemos:

\begin{equation}
f^{\prime}(n_t) = \frac{r\left( 1 + \frac{r - 1}{K} n_t \right) - r n_t \frac{r - 1}{K}}{\left( 1 + \frac{r - 1}{K} n_t \right)^2}
\label{eqn:derivada-beverton-holt}
\end{equation}

Al evaluar esta expresión en $n_t = K$ y $n_t = 0$, obtenemos que $f^{\prime}(K) = 1/r$ y $f^{\prime}(0) = r$. Por lo tanto, el punto de equilibrio $n_t = K$ es estable si $r > 1$ y el punto de equilibrio $n_t = 0$ es estable si $r < 1$.

Veamos esto ejemplificado en una simulación numérica con $K = 50$, realizada hasta un tiempo $t = 50$, donde se varían $r$ y la población inicial $n_0$.

\begin{itemize}
\item En la Figura \ref{fig:beverton-holt-r-0.5}, utilizamos un valor de $r = 0.5 < 1$. Se observa que el sistema tiende al punto de equilibrio $n_t = 0$, es decir, la población se extingue.
\item En la Figura \ref{fig:beverton-holt-r-2}, utilizamos un valor de $r = 2 > 1$,. Se observa que el sistema tiende al punto de equilibrio $n_t = K$, es decir, a la población máxima limitada por la capacidad de carga del modelo.
\item En la Figura \ref{fig:beverton-holt-r-1}, se observa que con $r = 1$ el sistema se queda en el punto inicial $n_0$, es decir, no hay crecimiento ni decrecimiento de la población.
\end{itemize}

\begin{figure}[H]
    \centering
    \includegraphics[width=1\linewidth]{../Practico/ej_01/Beverton_Holt_r_0.5.pdf}
    \caption{Simulación del mapeo de Beverton - Holt con $r = 0.5$ para distintas condiciones iniciales.}
    \label{fig:beverton-holt-r-0.5}
\end{figure}

\begin{figure}[H]
    \centering
    \includegraphics[width=1\linewidth]{../Practico/ej_01/Beverton_Holt_r_2.pdf}
    \caption{Simulación del mapeo de Beverton - Holt con $r = 2$ para distintas condiciones iniciales.}
    \label{fig:beverton-holt-r-2}
\end{figure}

\begin{figure}[H]
    \centering
    \includegraphics[width=1\linewidth]{../Practico/ej_01/Beverton_Holt_r_1.pdf}
    \caption{Simulación del mapeo de Beverton - Holt con $r = 1$ para distintas condiciones iniciales.}
    \label{fig:beverton-holt-r-1}
\end{figure}

\section{Ecuación logística con retraso}
La ecuación logística con retraso consiste en una variante de la ecuación logística que tiene en cuenta la historia previa del sistema hasta un tiempo $T$. Está definida por la siguiente ecuación diferencial:

\begin{equation}
\frac{dN}{dt} = rN(t)\left[ 1 - \frac{N(t - T)}{K} \right],
\label{eqn:logistic-delay}
\end{equation}

donde $r$ es la tasa de reproducción, $K$ es la capacidad de carga y $T$ es el retraso.

Se resolvió numéricamente esta ecuación variando la tasa de reproducción $r$, manteniendo $K = 10$ y $T = 1$ constantes. Dado que esta ecuación diferencial incluye un retraso, necesitamos condiciones iniciales para $t \in [-T, 0]$; en esta simulación, utilizamos $N(t) = 2$ para estas condiciones iniciales.

Los resultados obtenidos se muestran en la Figura \ref{fig:logistic-delay-num}. En ella se pueden observar distintos regímenes. Para $r = 0.3$, se observa un régimen monótono similar al de la ecuación logística sin retraso. Para $r = 1.2$, comienzan a observarse pequeñas oscilaciones amortiguadas antes de converger al valor de equilibrio. Mientras que para $r = 2$, se observa un régimen oscilatorio sin tender a un valor fijo.

Aunque en un sistema dinámico unidimensional no es posible que aparezcan oscilaciones, estas ocurren en este caso debido a que el sistema requiere de infinitas condiciones iniciales para su resolución.

\begin{figure}[H]
    \centering
    \includegraphics[width=1\linewidth]{../Practico/Ej_02/Logistica_delay.pdf}
    \caption{Simulación numérica de la solución de la ecuación logística con delay para distintas tasas de reproducción $r$.
    Pueden observarse distintos regímenes de la solución.}
    \label{fig:logistic-delay-num}
\end{figure}

Con el objetivo de verificar la solución numérica obtenida, se realizó una comparación con una solución aproximada de la ecuación (\ref{eqn:logistic-delay}), dada por:

\begin{equation}
N(t) \approx 1 + c e^{\frac{\epsilon t}{1 + \pi^2 / 4}} e^{i t \left[ 1 - \frac{\epsilon \pi}{2 \left( 1 + \pi^2/4 \right)} \right]},
\label{eqn:logistic-delay-aprox}
\end{equation}

Esta aproximación es válida para un régimen donde $T$ es mayor al valor crítico $T_c = \pi/2r$. En este caso, tomamos $T = \epsilon + T_c$, con $r = 2$ y un valor de $\epsilon = 1 \times 10^{-3}$ para las simulaciones.

Los resultados se muestran en la Figura \ref{fig:logistic-delay-aproxvnum}. Se observa que para tiempos cortos, las soluciones son aproximadamente iguales; sin embargo, para tiempos largos, las soluciones se desfasan. Además, se aprecia que la solución aproximada continúa sin decaimiento, mientras que esto sí ocurre para la solución numérica.

\begin{figure}[H]
    \includegraphics[width=1\linewidth]{../Practico/Ej_02/Aproximacion_vs_exacta.pdf}
    \caption{Solución numérica de la ecuación logística con delay comparada con la solución aproximada de la misma.}
    \label{fig:logistic-delay-aproxvnum}
\end{figure}

Para los regímenes con oscilaciones, se analizó la dependencia de la amplitud de las mismas con la condición inicial. Se realizaron simulaciones tomando $N_0 = {2,~6.5,~11}$ constantes para $t \in [-T, 0]$. Los resultados de las soluciones en estado estacionario se muestran en la Figura \ref{fig:vefirica_amplitudes}. Se observa que la amplitud de las oscilaciones no depende de la condición inicial. Esto también se verifica en la Figura \ref{fig:vefirica-dif-amplitudes}, donde se muestran las diferencias entre las amplitudes para distintas simulaciones, las cuales no superan una diferencia máxima de 0.07.

\begin{figure}[H]
    \includegraphics[width=1\linewidth]{../Practico/Ej_02/Verificacion_amplitud.pdf}
    \caption{Soluciones numéricas de la ecuación logística con delay con distintas condiciones iniciales con los parámetros para encontrarse en el régimen oscilatorio.}
    \label{fig:vefirica_amplitudes}
\end{figure}

\begin{figure}[H]
    \includegraphics[width=1\linewidth]{../Practico/Ej_02/Verificacion_diferencias_amplitud.pdf}
    \caption{Diferencias entre las amplitudes de distintas simulaciones en función de la condición inicial $N_0$.}
    \label{fig:vefirica-dif-amplitudes}
\end{figure}

Se realizó un análisis similar para investigar la dependencia del período de las oscilaciones en función de la tasa de reproducción $r$. Los resultados se muestran en la Figura \ref{fig:vefirica-periodo}. En esta figura se puede observar que el período aumenta a medida que se incrementa la tasa de reproducción $r$.
\begin{figure}[H]
    \includegraphics[width=1\linewidth]{../Practico/Ej_02/Verificacion_periodo.pdf}
    \caption{Período de las soluciones numéricas de la ecuación logística con delay en el régimen oscilatorio con distintos valores de $r$.}
    \label{fig:vefirica-periodo}
\end{figure}

\section{Modelo matricial de Leslie}
Las matrices de Leslie son un tipo de matriz utilizada en biología matemática para modelar la dinámica poblacional, especialmente en poblaciones estructuradas por edades. La solución
está asociada a un crecimiento exponencia único, $L E = r E$ donde $L$ es la matriz de Lesile para $k+1$ grupos etarios como:

\begin{align*}
& E_1=\frac{S_0}{r} E_0, \quad r  E_0=f_0 E_0+\ldots+f_n E_K \\
& E_2=\frac{S_1}{r} E_1=\frac{S_1 S_0}{r^2} E_{0}\\
& E_k=\frac{S_{k-1} \cdots S_0 E_0}{r^k} \\
\end{align*}

Con estas ecuaciones:
\begin{eqnarray}
    r^{k+1} - \left( r^k f_0 + r^{k-1} f_1 s_0 + r^{k-2} f_2 s_0 s_1 + \right. \nonumber \\
    \left. \ldots + f_k s_0 s_1 \cdots s_{k-1}\right).
    \label{eqn:leslie}
\end{eqnarray}
Una de las raíces es el valor de $r$ que buscamos, según el criterio de Decartes como hay un solo cambio de signo de los coeficientes 
tendremos una sola raíz positiva que llamaremos $r_1$

Evaluando $r_1$ en la ecuación \ref{eqn:leslie} obtenemos:

$$
R=1-\left(f_0 + f_1 s_0 + \ldots + f_k s_0 s_1 \cdots s_{k-1}\right)
$$

Supongamos que $r_1 > 1$. Al ser raíz del polinomio

\begin{align*}
& r_1^{k+1}-\left(r_1^k f_0 + r_1^{k-1} f_1 s_0 + \ldots + f_k s_0 s_1 \cdots s_{k-1}\right) = 0 \\
& r_1^{k + 1}\left[1-\frac{f_0}{r_1}+\frac{f_1 s_0}{r_1^2}+\cdots+\frac{f_k s_0 s_1 \cdots s_{k-1}}{r_1^{k+1}}\right]=0 \\
& \text { como } r_1 \neq 0 \\
& \frac{f_0}{r_1}+\frac{f_1 s_0}{r_1^2}+\cdots+\frac{f_k s_0 s_1 \cdots s_{k-1}}{r_1^{k+1}}=1 \\
& \Rightarrow \left(f_0 + f_1 s_0 + \ldots + f_k s_0 s_1 \cdots s_{k-1}\right) > 1,
\end{align*}

pues supusimos que $r_1 > 1$ entonces $R < 0$.

Ahora supongamos que $R < 0$ 

\begin{eqnarray*}
    1 - \left(f_0 + f_1 s_0+\ldots+f_k s_0 s_1 \ldots s_{k-1}\right) < 0 \\
    \Rightarrow 1<f_n+f_1 s_0+\ldots+f_k s_0 s_1 \ldots s_{k-1}
\end{eqnarray*}

Supongamos, que $r_1<1$ entonces

\begin{eqnarray*}
    \frac{f_0}{r_1}+\frac{f_1 s_0}{r_1^l}+\cdots+\frac{f_k s_0 s_1 \cdots s_{n-l}}{r_1^{k+1}} > \\
    f_0+f_1 s_0+\cdots+f_k s_0 s_1 \cdots s_{k-1}>1
\end{eqnarray*}

Pero vimos que si $r_1$ es raíz el término de la izquierda es 1 $\Rightarrow$ tenemos $1>1$, absurdo que viene de suponer $r_1<1$ y si $r_1 = 1$ entonces $R = 0$ lo que tampoco coincide con la suposición $R < 0$ por lo tanto $r_1 > 1$

\section{Especie con ciclo de vida anual}
Simulamos un sistema de individuos en el cual cada individuo produce $r$ descendientes y luego muere. La población evoluciona de acuerdo
a:
\begin{equation}
    N_{t+1} = r N_t
\end{equation}
Este $r$ obedece a una distribución de Poisson con media $1.7$ y comenzamos con un solo individuo, en la figura 
\ref{fig:poblacion_poisson} se muestra la evolución de la población en función del tiempo. 

En principio se observa que la población crece de manera exponencial, sin embargo dado que la descendencia muchas veces es 
0 puesto que $r$ es un valor aleatório tomado de una distribución de Poisson, ocurre que la población se termina extinguiendo.

\begin{figure}[H]
    \includegraphics[width=1\linewidth]{../Practico/Ej_04/SimulacionPoisson.pdf}
    \caption{Población $N$ en función del tiempo $t$ para una descendencia $r$ variable aleatoria que sigue una distribución de Poisson con media $1.7$.}
    \label{fig:poblacion_poisson}
\end{figure}

\section{Población de animales costero intertidal}

Considereamos una población de animales costeros, que viven entre las líneas de marea alta y baja. Estas poblaciones son vulerables al efecto de tormentas severas,
que las afectan de distinto modo según su intensidad y el estado de la marea. Supongamos que puede modelar el sistema mediante una evolución determinista entre 
eventos desastrosos que ocurren al azar.

Consideramos un modelo logistico entre desastres: $\dot{N} = rN(1-N/K)$, donde $r$ es la tasa de crecimiento y $K$ es la capacidad de carga. Si ocurre una castástrofe a tiempo $t$
la población se ve disminuida en una fracción $p$ como $N(t^{+}) = pN(t)$. Los tiempos entre desastres siguen una distribución exponencial con media $1/\lambda$.

Realizamos una simulación de este modelo durante el periodo $t \in [0,50]$ con $r = 1$, $K = 100$, $p = 0.5$, $\lambda = 0.5$ y población inicial $N_0 = 10$. 
En la figura \ref{fig:costero-intertidal} se muestra la evolución de la población en función del tiempo.

\begin{figure}[H]
    \includegraphics[width=1\linewidth]{../Practico/Ej05/CatastrofesPoblacion.pdf}
    \caption{Evolución de la población $N$ en función del tiempo $t$ para una población costera intertidal con eventos catastróficos.}
    \label{fig:costero-intertidal}
\end{figure}

Se puede observar como la población a pesar de sufirir eventos catastróficos, que se muestran como disminuciones abruptas en la población, esta logra recuperarse y llegar a un equilibrio estable.

Encontremos una condición para la extinción, sabemos que la ecuación logística entre catástrofes tiene una solución dada por

\begin{equation}
    N(t) = \frac{N_0 K e^{rt}}{K + N_0 \left( e^{rt} - 1 \right)},
    \label{eqn:SolLogistica}
\end{equation}

donde $N_0$ es la población inicial. Si suponemos que a un tiempo $t \approx \frac{1}{\lambda}$ ocurre una catástrofe entonces tendremos que en  un tiempo
posterior $t^{+}$

\begin{equation}
    N(t^{+}) = p N(\frac{1}{\lambda}) = \frac{p N_0 K e^{\frac{r}{\lambda}}}{K + N_0 \left( e^{\frac{r}{\lambda}} - 1 \right)}
    \label{eqn:SolLogisticaCatastrofe}
\end{equation}

Si continuamos con la evolución de la población depués de la catástrofe tenemos nuevamente la ecuación logística por lo que
la población a un tiempo $t_2$ será:

\begin{equation}
    N_2 (t_2) = \frac{N(t^{+}) K e^{r \left(t_2 - \frac{1}{\lambda} \right)}}{K + N(t^{+}) \left( e^{r \left(t_2 - \frac{1}{\lambda} \right)} - 1 \right)}
\end{equation}

Es decir nuevamente tenémos una solución de la logística pero esta vez con la condición inicial en $N(t^{+})$.
Si la población se extingue tenemos que considerar que $N(t^{+}) \approx 0$. Si hacemos esto en expansión de Taylor quedamos en que

\begin{equation}
    N(t) \approx p N(t^{+}) e^{\frac{r}{\lambda}},
    \label{eqn:CondExtincion}
\end{equation}

en donde ignoramos los índices y tomamos como si empezaramos a contar desde el tiempo de la catástrofe.
Si la población se extingue entonces el factor $p e^{\frac{r}{\lambda}}$ debe ser menor a 1.
En la simulación realizada teníamos $p = 0.5$ y $\lambda = 0.5$ por lo que la condición de extinción es $r \lessapprox 0.345$.

En la figura \ref{fig:costero-intertidal-extincion} se muestra la evolución de la población en función del tiempo para un valor de $r = 0.3$.

\begin{figure}[H]
    \includegraphics[width=1\linewidth]{../Practico/Ej05/CatastrofesPoblacionExtincion.pdf}
    \caption{Evolución de la población $N$ en función del tiempo $t$ para una población costera intertidal con eventos catastróficos y extinción.}
    \label{fig:costero-intertidal-extincion}
\end{figure}

\section{Efecto Alle}
El efecto Alle da cuenta de un fenómeno, descripto por W.C. Allee, asociado a la existencia de un número crítico mínimo de individuos para garantizar la supervivencia de la especie.
Se supone que a partir de cierto umbral, el tamaño poblacional es tan reducido que los individuos no se reproducen al no encontrarse con otros individuos de la misma población. 
Una manera de modelar este efecto es por medio de la siguiente ecuación:

\begin{equation}
    \frac{dN}{dt} = rN \left[ 1 - \frac{N}{K} \right] \left[ \frac{N}{A} - 1 \right]
    \label{eqn:Alle}
\end{equation}

Analizamos los equilibrios de la ecuación \ref{eqn:Alle} y la estabilidad de los mismos. 
Para ello calculamos los puntos de equilibrio de la ecuación \ref{eqn:Alle} y la derivada de la misma con respecto a $N$.

\begin{eqnarray*}
    f^{\prime} = r \left( 1 - \frac{N}{K} \right) \left( \frac{N}{A} - 1 \right) \\
    - ~ \frac{rN}{K} \left( \frac{N}{A} - 1 \right) \\
    + ~ \frac{rN}{A} \left( 1 - \frac{N}{K} \right)
\end{eqnarray*}

Evaluando en los puntos de equilibrio $N = 0$, $N = K$ y $N = A$ tenemos que $f^{\prime}(0) = -r$, $f^{\prime}(K) = r \left(1 - \frac{K}{A} \right)$ y $f^{\prime}(A) = r \left(1 - \frac{A}{K}\right)$.
De estas obtenemos las siguientes condiciones de equilibrio:

\begin{itemize}
    \item Si $r > 0$:
    \begin{itemize}
        \item $N = 0$ es estable
        \item $N = K$ es estable si $K > A$
        \item $N = A$ es estable si $A > K$
    \end{itemize}

    \item Si $r < 0$:
    \begin{itemize}
        \item $N = 0$ es inestable
        \item $N = K$ es estable si $K < A$
        \item $N = A$ es estable si $A < K$
    \end{itemize}

\end{itemize}

En la ecuación logística para una población buscabamos que $N = 0$ sea inestable para $r > 0$ y estable para $r < 0$. 
En este caso es distinto porque cuando $r>0$ la población se extingue si $A < K$, por lo tanto si empezamos con una condición inicial $N_0 < A$ la población se extingue.

Esto se muestra en la figura \ref{fig:alle} donde se simula la evolución de la población en función del tiempo 
para distintos valores de $N_0$ y valores fijos de $A = 10$, $K = 20$ y $r = 0.1$.

\begin{figure}[H]
    \includegraphics[width=1\linewidth]{../Practico/Ej06/Alle.pdf}
    \caption{Evolución de la población $N$ en función del tiempo $t$ para una población con efecto Alle.}
    \label{fig:alle}
\end{figure}

\end{multicols}

\end{document}